\section{模型的检验与灵敏度分析}
\label{sec:6}

\subsection{模型检验}
\label{sec:6.1}

全文检验汇总于表~\ref{tab:checks}。「互相验证」用于同一问题、不同假设的方法对，判据取排序或数值的一致性；上下游分工衔接的方法之间不构成验证关系。

\begin{table}[H]
\centering\footnotesize
\caption{全文检验清单}
\label{tab:checks}
\begin{tabularx}{\textwidth}{@{}c p{3.85cm}p{1.55cm}p{2.25cm}X@{}}
\toprule
问 & 检验对 & 类型 & 判据 & 结果 \\
\midrule
一 & Huber 主分 vs 均值/PCA/TOPSIS & 方法敏感性 & 样本 Spearman & $0.9936/0.8166/0.9588$ \\
一 & A1 vs A2/A3 & 扩展检验 & KS 的 $p$ 值 & $0.313/0.220$，未检出分布差异 \\
一 & 六来源候选跨集 & 规则复检 & 候选比例与方向 & $5.00\%/4.73\%/4.99\%$，主导方向有变化 \\
一 & clr+Huber vs 裸份额 OLS vs GBDT & 1M 留出 & 逐域/目标 $R^2$ & $0.772/0.260$、$0.654/0.418$、$0.923/0.870$ \\
一 & A1$\to$A18 质量桥接 & 五折＋六映射域 & $R^2$、MAE & $0.4106$、$0.0365$（$n=6$） \\
一 & 1M $\to$ 60M $\to$ 1B & 跨尺度（1B 独立） & 目标 Spearman & $0.523/0.467/0.677$ \\
一 & 尺度修正（幂律主口径） & 1B 独立检验 & 逐域 MAE & $2.953\to0.770$，降幅约 $74\%$ \\
一 & A12–A15 估算表 & 一致性对照 & 排序 & 与实测相反，作反例 \\
二 & 有地板 vs 无地板 & 嵌套 $F$ 检验 & $F$、$p$ & $F=307.11$，$p<10^{-15}$ \\
二 & $\kappa_N=\kappa_D$、$\kappa_D=0$ & 嵌套 $F$ 检验 & $F$、$p$ & $57.19$（$3.4\times10^{-13}$）拒绝；$0.63$（$0.43$）不拒绝 \\
二 & B7 新增 90 点 & 留出（插值） & RMSE、偏差 & $0.0531$、$+0.0219$；噪声代理 $0.0458$ \\
二 & 解析式 vs 数值解 & 互相验证 & 相对误差 & $<10^{-14}$（替代倍数）、$<10^{-8}$ \\
二 & $Q=1$ 切片 vs $L_0$ & 退化核对 & 残差区间 & $+0.0036$，$[-0.0120,0.0193]$ \\
二 & B4/B5/B2/B10 & 外部趋势对照 & Spearman & $0.983/0.959/0.788/1.000$，水平不可比 \\
二 & B11/B12 来源审计 & 元数据核对 & 覆盖、$D$ 记账 & 步数覆盖 $100\%$，中位比 $0.9986$ \\
三 & 数值解 vs 闭式 $N^*$ & 互相验证 & 相对误差 & $1.54\times10^{-7}$ \\
三 & 求解器 vs 粗网格 & 互相验证 & Loss 差 & 最大差 $1.80\times10^{-5}$ \\
三 & 剖面导数 & 有限差分 & 相对误差 & $4.1\times10^{-10}$ \\
四 & C8 重算 vs C1 榜单 & 口径核验 & Pearson、MAE & $0.9893$–$0.9995$；MATH $0.8009/3.4974$ \\
四 & S 形 vs 线性 vs 常数 & 留一 & RMSE & $8.079/8.283/9.514$ \\
四 & 分解三口径 & 方法敏感性 & 技术占比极差 & $32.81$ pp \\
四 & 预测回测 & 事后回测 & 2025.25 误差 & M1b $-0.784$；朴素线性 $+2.651$ \\
\bottomrule
\end{tabularx}
\end{table}

\textbf{问题一。} 检验分三层：A2/A3 为冻结评分的扩展检验，A6/A7 为配比模型的同尺度留出，A8–A11 为跨尺度检验，其中 A10/A11（1B）完全不参与估计、A8/A9 承担尺度修正标定。GBDT 依据 A6/A7 表现选为交叉评估器，该留出不再重复充当候选收益的独立验证。A12–A15 为估算/外推表，其隐含的大尺度排序反转与实测排序保持相矛盾，作反例使用。

\textbf{问题二。} 质量参数只用 B6（半合成）估计，留出只用 B7 新增 90 点（$Q$ 方向插值）。B2 为半合成、B3 由 B1 插值派生、B10 为公式估算，二者用于核对实现与趋势；B4 的 57 点剔除 13 个近邻后 44 点结果稳定（Spearman $0.975$）。B8 与 B6/B7 方向相反（$64.26\%$ 低于 $E$），作反例。B11/B12 审计确认 B1 的 8 个名义规模、每模型 147 步均被覆盖，B1 训练设定真实、损失按公开公式生成。

\textbf{问题三。} 数值最优解与闭式 $N^*$、求解器与独立粗网格、剖面解析导数与有限差分三组互验，相对误差分别为 $1.54\times10^{-7}$、最大 Loss 差 $1.80\times10^{-5}$、$4.1\times10^{-10}$，预算残差与份额和误差均小于 $10^{-9}$；锚点映射与单域越界均按规范标记。

\textbf{问题四。} C8 重算与 C1 为同批数据的口径核验（MATH 差异待解释），桥接为同域留一，分解差异为方法敏感性，回测为事后回测且重估样本未严格截在起点前。问题四目前没有独立的时间外检验层，结论归属拟合层与条件外推层，这一边界在 \ref{sec:7.2} 说明。

\subsection{灵敏度与稳健性分析}
\label{sec:6.2}

\textbf{问题一。} 六来源阈值取 P90/P95/P97.5 时候选比例约 $10.00\%/5.00\%/2.50\%$（A2 $9.19\%/4.73\%/2.00\%$，A3 $9.74\%/4.99\%/2.44\%$），阈值直接决定人工复核工作量。折扣对照分 $Q_{\rm v2}$ 与主分秩相关达 $0.9547/0.9207/0.9395$、域级排序不变。单指标极端扰动下主分变动约小 $30\%$；13 域等权与 pile\_cc 单域两目标给出不同候选，评价目标在搜索前固定；质量交互项仅 6/13 域显著。

\textbf{问题二。} 去地板后质量 $+0.1$ 在千亿参数处的等效倍数由约 $3.0$ 降到 $1.3$–$1.5$，方向不变；$\kappa_N$、$k_0$ 的 $95\%$ 区间为 $[0.1264,0.1568]$、$[0.2013,0.2554]$，$\kappa_D$ 区间含 0；主式形状代价为留出 RMSE $0.0531$（线性对照 $0.0416$、噪声 $0.0458$）。$C=10^{19}$ 时提质是否划算取决于成本函数，$C\ge10^{20}$ 后三种 $g(Q)$ 下提质都更有效。

\textbf{问题三。} 敏感性覆盖三类成本、C7 五档上下文、$\eta$ 三档、P2 重抽样、两种质量映射与配比不确定性，以及 \texttt{m=1}、H$_{\rm tot}$、无地板对照。与预算和规模无关的解析结论是 $L_{ctx}^{\rm crit}=30000$ token；仅含 P2 时，离开下界的 bootstrap 区间为指数型 $[19.038,19.084]$、幂函数型 $[19.503,19.560]$、对数型 $[18.539,18.620]$。Q3-9 已将 P1 域质量与配比 bootstrap 与 P2 联合传播，离开下界区间外扩约 $0.003$ dex，三种成本函数的阈值排序不变。

\textbf{问题四。} Medium 层权重在 $0.25$–$1.0$ 时 $S(1.85)$ 变化不足 $0.4$ 分，仅用 High 层时上平台降到 $6.349$，上升段依赖 Medium 层。C3 权重取 $0.2$、$1.0$ 时技术率为 $0.499$、$0.205$ dex/年（不含 C3 为 $0.781$）。换窗口后 M1b 技术占比为 $7.06\%$、$8.18\%$，「规模为主」稳定。预测对后训练增益是否延续最敏感（约 $6.8$ 分），大于算力情景差异（约 $2$ 分）。

\section{模型的评价与推广}
\label{sec:7}

\subsection{模型的优点}
\label{sec:7.1}

\begin{enumerate}
\item \textbf{唯一质量口径与清晰证据分层。} 文档、领域、桥接标签与配比加权质量统一使用 CRITIC 加权 Huber 主分；A1 拟合、A2/A3 扩展、A6/A7 留出、A8/A9 标定、A10/A11 独立检验、A12–A15 估算分层报告。
\item \textbf{训练与检验边界可追溯。} 问题二只用 B6 估计质量参数、选模型，B7 新增 90 点作留出；配比通道参数只由 1M 训练行确定，A6–A11 仅用于评估。
\item \textbf{形式与阈值可检验。} 六来源阈值取 A1 同域 P95，桥接以分层五折比较 Ridge 与 GBDT；结构性转移由质量约束活跃状态定义并区分局部根与全局切换；主式显式指定并同时报告交叉验证更优的线性对照。
\item \textbf{解析与数值互验。} 替代条件、$N^*$ 与算力等价倍数、降维闭式与 $\Phi$ 判据、S 形桥接与 Shapley 分解均有闭式或可核验形式，并与数值结果逐项对照。
\item \textbf{职责分离、避免循环论证。} 先固定主分再报告分歧诊断与旧法对照，阈值敏感性与折扣对照分均不回写主评分，筛查预算不被当作真实冲突率。
\end{enumerate}

\subsection{模型的局限与改进方向}
\label{sec:7.2}

本研究的证据边界集中说明如下，凡正文结论涉及这些条件时均以此为准。

\begin{enumerate}
\item \textbf{无独立质量真值。} 问题一以方法一致性、扰动稳健性与下游敏感性刻画质量口径，不给出质量准确率；17 域中 11 域质量依赖文本桥接（五折 $R^2=0.4106$），区间不含模型选择与跨域偏差，ubuntu\_irc（16 篇）、philpapers（63 篇）低置信使用。
\item \textbf{质量参数来自半合成数据。} $\kappa_N,\kappa_D,k_0$ 只由 B6（半合成，最大约 12B）估计，B7 留出为 $Q$ 方向插值；主式为指定幂次，留出 RMSE（$0.0531$）高于线性对照（$0.0416$）与噪声代理（$0.0458$），$\kappa_D$ 较弱；高预算端地板份额（约 $33\%$）与「提质远比扩参数划算」含地板外推成分。
\item \textbf{跨尺度外推幅度不足。} 1B 检验显示尺度修正把绝对误差降低约 $74\%$，但平均预测仍高于观测，配比效应幅度在 1B 缩至约六成；两档规模下单参数缩放（对数线性或幂律）刻画不出曲率，曲率参数欠定。
\item \textbf{问题三的数据边界。} Q3-9 已闭合固定配比与重选配比两类抽样；质量参数仍来自半合成 B6，配比通道为探索性，$\log_{10}C\gtrsim23.5$ 的最优规模配置为外推。
\item \textbf{问题四口径尚未由重跑固定。} 主面板未执行严格开放权重纳入与许可证审计，C6 桥接迁移误差未单独进入预测区间，预测以历史锚点 $T_0=2025.25$ 为起点且原面板含前视数据，主链未在路径修复后重跑，补充独立实现尚未在本仓库运行（\texttt{DEFERRED\_Q3}）。
\item \textbf{问题四方法分歧与占比口径。} C6 的 High 层仅 7 点且全在底噪区，上升段依赖 Medium 层；分解三口径技术占比极差达 $32.81$ pp，是最大不确定性来源；占比分母为模型解释增量而非观测增量，观测与模型增量不完全闭合。
\item \textbf{预测不含离散突破。} 条件预测不含新架构等非平滑变化；个别低算力情景区间下界低于锚点，与累计前沿不下降定义存在张力，需先修复单调性。
\end{enumerate}

\textbf{改进方向。}
\begin{enumerate}
\item 为 11 个 inferred 域增补质量标签或模型评分信号，重新校准并检验桥接误差；补齐 P1 逐次重抽样接口，使两类配比不确定性可分别传播到问题三。
\item 取得真实跨质量、跨规模训练日志后替换 B6 重估 $k_0$ 与 $\kappa_N$，并允许配比强度随 $\log N$ 变化；在获得真正跨尺度配比试验前对 H$_{\rm red}$、H$_{\rm tot}$ 各算一遍。
\item 问题四重跑时执行明确的开放权重纳入规则（报告纳入前后样本量与敏感性），把建模、校准与回测样本全部截在各自起点之前，将 C6 桥接的 bootstrap 与残差作为独立不确定性通道，并统一占比分母与前沿单调性。
\item 取得更多同验证集的 Loss–Benchmark 对，缩小结构口径与线性口径的技术占比分歧。
\end{enumerate}

\subsection{模型的推广}
\label{sec:7.3}

「经典基线＋质量修正＋配比乘子」的标度律结构可推广到其他数据受限的资源配置问题（去重、合成数据、多模态对齐均可写成可约损失上的乘子与不随规模消失的地板）；「约束取等后解析消元＋活跃约束状态分类」的降维方法适用于含多种数据处理成本的预算优化；「分层 S 形桥接＋连续时间项＋Shapley 分解」框架适用于其他存在能力饱和的技术演进分析。推广时保持同一纪律：明确数据证据等级，把半合成、估算与外推数据标注为假设支撑而非观测验证。

\section{结论}
\label{sec:8}

本文围绕「算力约束下如何配置资源以提升大语言模型能力」，建立四个层层衔接的模型，主要结论如下。

\textbf{问题一。} 以 A1 冻结的 CRITIC 加权 Huber 中心作为唯一质量口径，把 22 个字段展开为 25 项指标，得到文档、领域与语料质量；六来源诊断在三集各圈定约 $5\%$ 的高分歧复核候选，主导对立方向随语料变化；17 域质量由映射与文本桥接补齐（五折 $R^2=0.4106$）。配比模型给出 $17\times13$ 迁移结构，推荐配比 $\bp^*$ 在 1M 代理模型下较均匀配比预测降损 $9.47\%$（$95\%$ 区间 $[7.82\%,11.16\%]$）；域质量与训练边际价值秩相关仅 $-0.211$，质量分不能替代配比实验。

\textbf{问题二。} 建立幂次广义标度律 $L=E+k_0(1-Q)+[AN^{-\alpha}Q^{-\kappa_N}+BD^{-\beta}Q^{-\kappa_D}]\exp(s_{\rm red}\tilde\varphi(\bp))$，其中经典参数 $E=1.6901$、$A=406.39$、$\alpha=0.34$、$B=410.59$、$\beta=0.28$，质量参数 $\kappa_N=0.1416$、$\kappa_D=0.0112$、$k_0=0.2280$、$Q_0=0.6624117$。参照点上质量弹性约为参数弹性的 $1.62$ 倍；质量提高 $0.1$ 等效于参数扩大 $1.30$–$3.07$ 倍；沿算力前沿等效算力乘 $1.191/1.485/2.134$。$C\ge10^{20}$ 后提质比扩参数更有效；配比通道受单一尺度限制，按探索性下、上情景报告。

\textbf{问题三。} 约束取等后解析消元，把问题降为 $(\log_{10}N,Q)$ 二维有界优化，以质量约束活跃状态定义结构性转移，以 $\Phi=1$ 作为固定 $Q$ 下的局部一阶必要条件（导数核验误差约 $4.1\times10^{-10}$）。上下文长度临界值为 $L_{ctx}^{\rm crit}=30000$ token。数值上，$C=10^{19}$ 时指数型、幂函数型成本下不提质、对数型直接提满，$C=10^{22},10^{24}$ 时均提满；离开下界的临界预算（$\log_{10}C$）约为 $19.06/19.53/18.58$，为「何时投向提质」给出随成本函数移动的明确阈值。

\textbf{问题四。} 用「标度律＋S 形桥接＋连续时间项」分解历史前沿：$2023.50\to2025.00$ 窗口 base 观测提高 $24.759$ 分，规模贡献 $23.063$、技术贡献 $3.404$，技术占比 $12.86\%$（$90\%$ 区间 $6.07\%$–$21.96\%$）；去饱和综合分同期提高 $48.19\%$，更灵敏地反映进展。算力增速由 $0.60$ 降至 $0.15$ dex/年时，24 个月 base 中位数为 $44.11$，技术增量占比由 $54.94\%$ 升至 $76.70\%$。该结果以 $T_0=2025.25$ 为历史锚点，更新数据后需重作。

四问共同指向核心判断：\textbf{数据质量与领域配比是与算力并列的一等资源}。低预算阶段边际投入主要投向模型规模，预算跨过质量收益与提质开销的平衡点后，提质与配比优化优先级上升；算力增长放缓时，能力进步将越来越依赖非规模技术创新。上述定量判断的适用条件已在 \ref{sec:7.2} 逐条界定。

\begin{thebibliography}{99}\small
\bibitem{hoffmann} Hoffmann J, Borgeaud S, Mensch A, et al. Training compute-optimal large language models[C]//Advances in Neural Information Processing Systems 35 (NeurIPS). 2022.
\bibitem{bahri} Bahri Y, Dyer E, Kaplan J, et al. Explaining neural scaling laws[J]. Proceedings of the National Academy of Sciences, 2024, 121(27): e2311878121.
\bibitem{critic} Diakoulaki D, Mavrotas G, Papayannakis L. Determining objective weights in multiple criteria problems: The CRITIC method[J]. Computers \& Operations Research, 1995, 22(7): 763–770.
\bibitem{huber} Huber P J. Robust estimation of a location parameter[J]. The Annals of Mathematical Statistics, 1964, 35(1): 73–101.
\bibitem{fineweb} Penedo G, et al. The FineWeb datasets: Decanting the web for the finest text data at scale[C]//NeurIPS Datasets and Benchmarks Track. 2024.
\bibitem{aitchison} Aitchison J. The Statistical Analysis of Compositional Data[M]. London: Chapman \& Hall, 1986.
\bibitem{regmix} Xia M, et al. RegMix: Data mixture as regression for language model pre-training[C]//International Conference on Machine Learning (ICML). 2024.
\bibitem{pythia} Biderman S, et al. Pythia: A suite for analyzing large language models across training and scaling[C]//International Conference on Machine Learning (ICML). 2023.
\bibitem{muennighoff} Muennighoff N, et al. Scaling data-constrained language models[C]//Advances in Neural Information Processing Systems 36 (NeurIPS). 2023.
\bibitem{schaeffer} Schaeffer R, Miranda B, Koyejo S. Are emergent abilities of large language models a mirage?[C]//Advances in Neural Information Processing Systems 36 (NeurIPS). 2023.
\bibitem{shapley} Shapley L S. A value for $n$-person games[M]//Contributions to the Theory of Games II. Princeton: Princeton University Press, 1953: 307–317.
\bibitem{shafer} Shafer G. A Mathematical Theory of Evidence[M]. Princeton: Princeton University Press, 1976.
\bibitem{subramanyam} Subramanyam A, Chen Y, Grossman R L. Scaling laws revisited: Modeling the role of data quality in language model pretraining[J]. arXiv preprint arXiv:2510.03313, 2025.
\bibitem{softq} Xu Z, Wu S, Cho H, et al. Data-constrained language model pretraining: Improved regularization and scaling laws[J]. arXiv preprint arXiv:2606.06888, 2026.
\bibitem{lovelace} Lovelace J, Belardi C, Kundurthy S, et al. Prescriptive scaling laws for data constrained training[J]. arXiv preprint arXiv:2605.01640, 2026.
\end{thebibliography}
